% Generated mechanically from frozen input p12/ko/fdl.tex.
% Do not edit this generated file; edit the bound locale source instead.

% OK, start here.
%
%---------------------------------------------------------------------
%\tableofcontents

\setcounter{chapter}{115}
\chapter{GNU 자유 문서 사용 허가서 (GNU Free Documentation License)}

\phantomsection
\label{fdl-section-phantom}

\begin{center}

      버전 1.2, 2002년 11월


저작권 \copyright 2000, 2001, 2002  Free Software Foundation, Inc.

\bigskip

    51 Franklin St, Fifth Floor, Boston, MA  02110-1301  USA

\bigskip

누구든지 이 사용 허가서 문서의 원문 그대로의 사본을 복제하고
배포할 수 있으나, 이를 변경하는 것은 허용되지 않습니다.

비공식 번역 고지: 이것은 GNU 자유 문서 사용 허가서를 한국어로
옮긴 비공식 참고 번역입니다. 이 번역은 자유 소프트웨어 재단이 발행한
것이 아니며, GNU FDL을 사용하는 문서의 배포 조건을 법적으로 명시하지
않습니다. 그러한 효력은 GNU FDL의 원문 영어판에만 있습니다. 그러나 이
번역이 한국어 사용자가 GNU FDL을 더 잘 이해하는 데 도움이 되기를
바랍니다.

This is an unofficial translation of the GNU Free Documentation
License into Korean. It was not published by the Free Software Foundation,
and does not legally state the distribution terms for documentation that
uses the GNU FDL---only the original English text of the GNU FDL does that.
However, we hope that this translation will help Korean speakers understand
the GNU FDL better.

귀하는 수정 여부와 관계없이
http://www.gnu.org/licenses/translations.html 의 조건에 따라서만 이 번역을
게시할 수 있습니다.

You may publish this translation, modified or unmodified, only under the
terms at http://www.gnu.org/licenses/translations.html.
\end{center}


\begin{center}
{\bf\large 전문 (Preamble)}
\end{center}

이 사용 허가서의 목적은 설명서, 교과서 또는 그 밖의 기능적이고
유용한 문서를 자유라는 의미에서 "자유롭게" 만드는 것입니다. 즉,
상업적이든 비상업적이든, 수정하든 수정하지 않든, 누구에게나 그
문서를 복제하고 재배포할 실질적인 자유를 보장하는 것입니다.
부차적으로, 이 사용 허가서는 저자와 발행자가 자신의 작업에 대한
공로를 인정받을 수단을 보전하면서도, 다른 사람이 가한 수정에 대해
책임이 있는 것으로 간주되지 않게 합니다.

이 사용 허가서는 일종의 "카피레프트"입니다. 이는 문서의 파생물도
그와 같은 의미에서 자유로워야 한다는 뜻입니다. 이 사용 허가서는
자유 소프트웨어를 위해 마련된 카피레프트 사용 허가서인 GNU 일반
공중 사용 허가서(GNU General Public License)를 보완합니다.

자유 소프트웨어에는 자유 문서가 필요하므로, 우리는 자유 소프트웨어의
설명서에 사용하기 위해 이 사용 허가서를 설계했습니다. 자유 프로그램에는
그 소프트웨어와 같은 자유를 제공하는 설명서가 함께 있어야 합니다.
그러나 이 사용 허가서는 소프트웨어 설명서에만 한정되지 않으며, 주제나
인쇄 도서로 출판되는지 여부와 관계없이 모든 텍스트 저작물에 사용할 수
있습니다. 특히 교육이나 참고를 목적으로 하는 저작물에 이 사용 허가서를
사용할 것을 권장합니다.


\section{적용 범위와 정의 (APPLICABILITY AND DEFINITIONS)}
\label{fdl-section-applicability-and-definitions}

이 사용 허가서는 매체를 불문하고, 저작권자가 이 사용 허가서의 조건에
따라 배포할 수 있다고 명시한 고지가 포함된 모든 설명서나 그 밖의
저작물에 적용됩니다. 그러한 고지는 여기에 명시된 조건에 따라 그
저작물을 사용할 수 있는, 전 세계에 적용되고 사용료가 없으며 존속
기간에 제한이 없는 사용 허가를 부여합니다. 아래에서 \textbf{"문서"
("Document")}란 그러한 설명서나 저작물을 말합니다. 일반 공중의 모든
구성원은 사용 허가를 받은 자이며 \textbf{"귀하" ("you")}라고 지칭됩니다.
저작권법상 허가가 필요한 방식으로 저작물을 복제, 수정 또는 배포하면
귀하는 이 사용 허가를 수락한 것입니다.

문서의 \textbf{"수정판" ("Modified Version")}이란 문서 또는 그 일부를
원문 그대로 복제하거나, 수정하거나 다른 언어로 번역하여 포함한 모든
저작물을 말합니다.

\textbf{"부차적 절" ("Secondary Section")}이란 문서의 이름 붙은 부록 또는
앞부분의 절로서, 문서의 발행자나 저자와 문서의 전체 주제(또는 관련
사안) 사이의 관계만을 다루고 그 전체 주제에 직접 해당할 수 있는 내용은
전혀 포함하지 않는 것을 말합니다. (따라서 문서가 일부라도 수학
교과서라면, 부차적 절에서 어떠한 수학 내용도 설명해서는 안 됩니다.)
그 관계는 해당 주제 또는 관련 사안과의 역사적 연관일 수도 있고, 그에
관한 법적, 상업적, 철학적, 윤리적 또는 정치적 입장일 수도 있습니다.

\textbf{"변경 불가 절" ("Invariant Sections")}이란 문서가 이 사용 허가서에
따라 공개된다고 밝히는 고지에서 변경 불가 절이라고 그 제목을 지정한
특정 부차적 절을 말합니다. 어떤 절이 위의 부차적 절 정의에 부합하지
않으면 변경 불가 절로 지정할 수 없습니다. 문서에는 변경 불가 절이
하나도 없을 수 있습니다. 문서가 변경 불가 절을 식별하지 않으면 변경
불가 절은 없는 것입니다.

\textbf{"표지 문구" ("Cover Texts")}란 문서가 이 사용 허가서에 따라
공개된다고 밝히는 고지에서 앞표지 문구(Front-Cover Texts) 또는 뒷표지
문구(Back-Cover Texts)로 열거된 특정한 짧은 문구를 말합니다. 앞표지
문구는 최대 5단어, 뒷표지 문구는 최대 25단어일 수 있습니다.

문서의 \textbf{"투명" ("Transparent")} 사본이란 그 형식의 명세가 일반
공중에게 공개되어 있는 기계 판독 가능 사본으로서, 일반적인 텍스트
편집기, (픽셀로 구성된 이미지의 경우) 일반적인 그림 편집 프로그램,
또는 (도면의 경우) 널리 구할 수 있는 도면 편집기로 문서를 어렵지 않게
수정하기에 적합하고, 텍스트 조판 프로그램에 입력하거나 텍스트 조판
프로그램에 입력하기 적합한 여러 형식으로 자동 변환하기에 적합한 사본을
말합니다. 다른 면에서는 투명한 파일 형식으로 만들어졌더라도, 마크업이나
마크업의 부재가 독자의 후속 수정을 방해하거나 단념시키도록 구성된
사본은 투명하지 않습니다. 상당한 분량의 텍스트에 사용된 이미지 형식은
투명하지 않습니다. "투명"하지 않은 사본을 \textbf{"불투명" ("Opaque")}
사본이라고 합니다.

투명 사본에 적합한 형식의 예에는 마크업이 없는 일반 ASCII, Texinfo
입력 형식, LaTeX 입력 형식, 공개적으로 이용 가능한 DTD를 사용하는 SGML
또는 XML, 그리고 표준을 준수하며 사람이 수정할 수 있도록 설계된 단순한
HTML, PostScript 또는 PDF가 포함됩니다. 투명한 이미지 형식의 예에는
PNG, XCF 및 JPG가 포함됩니다. 불투명 형식에는 독점 워드 프로세서로만
읽고 편집할 수 있는 독점 형식, DTD 및/또는 처리 도구를 일반적으로 이용할
수 없는 SGML 또는 XML, 그리고 일부 워드 프로세서가 출력만을 목적으로
생성한 기계 생성 HTML, PostScript 또는 PDF가 포함됩니다.

\textbf{"제목 페이지" ("Title Page")}란 인쇄 도서의 경우 제목 페이지
자체와, 이 사용 허가서가 제목 페이지에 표시하도록 요구하는 내용을
읽을 수 있게 담는 데 필요한 뒤이은 페이지들을 말합니다. 제목 페이지가
따로 없는 형식의 저작물에서 "제목 페이지"란 본문이 시작되기 전에
나오는, 저작물의 제목이 가장 두드러지게 표시된 곳 근처의 텍스트를
말합니다.

\textbf{"XYZ라는 제목의" ("Entitled XYZ")} 절이란 그 제목이 정확히
XYZ이거나, XYZ를 다른 언어로 번역한 텍스트 뒤의 괄호 안에 XYZ가 들어
있는 문서의 이름 붙은 하위 단위를 말합니다. (여기에서 XYZ는 아래에서
언급하는 특정 절 이름, 예를 들어 \textbf{"감사의 글" ("Acknowledgements")},
\textbf{"헌사" ("Dedications")}, \textbf{"추천사" ("Endorsements")}, 또는
\textbf{"이력" ("History")}을 뜻합니다.) 귀하가 문서를 수정할 때 그러한
절의 \textbf{"제목을 유지" ("Preserve the Title")}한다는 것은 이 정의에
따라 그 절이 계속해서 "XYZ라는 제목의" 절로 남는다는 뜻입니다.

문서에는 이 사용 허가서가 문서에 적용된다고 밝히는 고지 옆에 보증 부인
조항(Warranty Disclaimers)이 포함될 수 있습니다. 이 보증 부인 조항은
참조에 의해 이 사용 허가서에 포함된 것으로 간주되지만, 보증을 부인하는
범위에 한합니다. 이 보증 부인 조항이 가질 수 있는 그 밖의 함의는 모두
무효이며 이 사용 허가서의 의미에 아무런 영향도 미치지 않습니다.


\section{원문 그대로의 복제 (VERBATIM COPYING)}
\label{fdl-section-verbatim-copying}

귀하는 상업적이든 비상업적이든 모든 매체로 문서를 복제하고 배포할 수
있습니다. 다만 이 사용 허가서, 저작권 고지 및 이 사용 허가서가 문서에
적용된다고 밝히는 사용 허가 고지를 모든 사본에 복제해야 하며, 이 사용
허가서의 조건에 어떠한 다른 조건도 추가해서는 안 됩니다. 귀하는 자신이
만들거나 배포하는 사본을 읽거나 다시 복제하는 일을 방해하거나 통제하기
위해 기술적 조치를 사용할 수 없습니다. 다만 사본을 제공하는 대가로
보상을 받을 수 있습니다. 충분히 많은 수의 사본을 배포하는 경우에는
제3절의 조건도 따라야 합니다.

귀하는 위에 명시된 것과 같은 조건으로 사본을 대여할 수도 있고,
사본을 공개적으로 전시할 수도 있습니다.


\section{대량 복제 (COPYING IN QUANTITY)}
\label{fdl-section-copying-in-quantity}


문서의 인쇄 사본(또는 통상 인쇄 표지가 있는 매체의 사본)을 100부 넘게
출판하고 문서의 사용 허가 고지가 표지 문구를 요구하는 경우, 모든 표지
문구를 명확하고 읽을 수 있게 담은 표지를 사본에 씌워야 합니다. 앞표지에는
앞표지 문구를, 뒷표지에는 뒷표지 문구를 실어야 합니다. 양쪽 표지에는
또한 귀하가 그 사본들의 발행자임을 명확하고 읽을 수 있게 표시해야
합니다. 앞표지에는 제목의 모든 단어가 똑같이 두드러지고 눈에 띄도록
완전한 제목을 표시해야 합니다. 표지에 그 밖의 자료를 추가할 수 있습니다.
변경이 표지에만 한정되고 문서의 제목을 유지하며 이 조건들을 충족하는
복제는 다른 면에서는 원문 그대로의 복제로 취급할 수 있습니다.

어느 한쪽 표지에 필요한 문구가 너무 길어서 읽을 수 있게 들어가지 않는
경우, 열거된 문구 가운데 앞의 것들부터 (합리적으로 들어가는 만큼)
실제 표지에 넣고 나머지는 인접한 페이지에 이어서 넣는 것이 좋습니다.

문서의 불투명 사본을 100부 넘게 출판하거나 배포하는 경우, 각 불투명
사본과 함께 기계 판독 가능한 투명 사본을 포함하거나, 각 불투명 사본
안이나 그와 함께 컴퓨터 네트워크상의 위치를 명시해야 합니다. 이 위치는
일반적으로 네트워크를 이용하는 공중이 공개 표준 네트워크 프로토콜로
문서의 완전한 투명 사본을 부가 자료 없이 내려받을 수 있는 곳이어야
합니다. 후자의 방식을 택하는 경우, 불투명 사본의 대량 배포를 시작할 때
합리적으로 신중한 조치를 취하여, 귀하가 그 판의 불투명 사본을 (직접
또는 대리인이나 소매상을 통하여) 공중에게 마지막으로 배포한 때로부터
적어도 1년 동안 그 투명 사본이 명시된 위치에서 계속 이와 같이 접근
가능하도록 보장해야 합니다.

의무 사항은 아니지만, 많은 수의 사본을 재배포하기 훨씬 전에 문서의
저자에게 연락하여 그들이 문서의 갱신판을 제공할 기회를 주기를
요청합니다.


\section{수정 (MODIFICATIONS)}
\label{fdl-section-modifications}

귀하는 위 제2절 및 제3절의 조건에 따라 문서의 수정판을 복제하고 배포할
수 있습니다. 다만 수정판을 정확히 이 사용 허가서에 따라 공개하고,
수정판이 문서의 역할을 하게 하여, 그 사본을 소유한 누구에게나 수정판의
배포와 수정을 허가해야 합니다. 또한 수정판에서 다음 사항을 이행해야
합니다.

\begin{itemize}
\item[A.]
   제목 페이지(그리고 표지가 있다면 표지)에 문서의 제목 및 이전 판들의
   제목과 구별되는 제목을 사용하십시오. (이전 판이 있었다면 그 제목들은
   문서의 이력 절에 열거되어 있는 것이 좋습니다.) 이전 판의 원 발행자가
   허가하면 그 이전 판과 같은 제목을 사용할 수 있습니다.

\item[B.]
   제목 페이지에 수정판의 수정 내용에 대한 저작에 책임이 있는 한 명
   이상의 개인 또는 단체를 저자로 열거하고, 아울러 문서의 주요 저자 중
   적어도 다섯 명(주요 저자가 다섯 명보다 적으면 그 전원)을 열거하십시오.
   다만 그들이 귀하에게 이 요구 사항을 면제해 준 경우는 예외입니다.

\item[C.]
   제목 페이지에 수정판의 발행자 이름을 발행자로서 명시하십시오.

\item[D.]
   문서의 모든 저작권 고지를 보존하십시오.

\item[E.]
   다른 저작권 고지에 인접하여 귀하의 수정에 대한 적절한 저작권 고지를
   추가하십시오.

\item[F.]
   저작권 고지 바로 뒤에, 아래 부록에 제시된 형식으로 이 사용 허가서의
   조건에 따라 공중이 수정판을 사용할 수 있도록 허가하는 사용 허가
   고지를 포함하십시오.

\item[G.]
   그 사용 허가 고지에서 문서의 사용 허가 고지에 제시된 변경 불가 절과
   필수 표지 문구의 전체 목록을 보존하십시오.

\item[H.]
   이 사용 허가서의 변경되지 않은 사본을 포함하십시오.

\item[I.]
   "이력" ("History")이라는 제목의 절을 보존하고 그 제목을 유지하며,
   제목 페이지에 제시된 수정판의 제목, 연도, 새 저자 및 발행자를 적어도
   명시하는 항목을 그 절에 추가하십시오. 문서에 "이력"이라는 제목의 절이
   없다면, 제목 페이지에 제시된 문서의 제목, 연도, 저자 및 발행자를
   명시하는 절을 새로 만든 다음, 앞 문장에서 정한 대로 수정판을 설명하는
   항목을 추가하십시오.

\item[J.]
   문서의 투명 사본에 공중이 접근하도록 문서에 제시된 네트워크 위치가
   있다면 이를 보존하고, 마찬가지로 문서의 토대가 된 이전 판들에 대해
   문서에 제시된 네트워크 위치도 보존하십시오. 이 위치들은 "이력" 절에
   둘 수 있습니다. 문서 자체보다 적어도 4년 전에 출판된 저작물의
   네트워크 위치는 생략할 수 있으며, 그 위치가 가리키는 판의 원 발행자가
   허가한 경우에도 생략할 수 있습니다.

\item[K.]
   "감사의 글" ("Acknowledgements") 또는 "헌사" ("Dedications")라는
   제목의 모든 절에 대해 그 절의 제목을 유지하고, 그 절에 제시된 각
   기여자에 대한 감사 및/또는 헌사의 모든 내용과 어조를 보존하십시오.

\item[L.]
   문서의 모든 변경 불가 절을 그 본문과 제목을 변경하지 않은 채
   보존하십시오. 절 번호나 그에 상응하는 표시는 절 제목의 일부로
   간주되지 않습니다.

\item[M.]
   "추천사" ("Endorsements")라는 제목의 모든 절을 삭제하십시오. 그러한
   절은 수정판에 포함할 수 없습니다.

\item[N.]
   기존 절의 제목을 바꾸어 "추천사"라는 제목을 붙이거나 변경 불가 절의
   제목과 충돌하게 해서는 안 됩니다.

\item[O.]
   모든 보증 부인 조항을 보존하십시오.
\end{itemize}

수정판에 부차적 절의 요건을 충족하고 문서에서 복제한 자료를 포함하지
않는 새로운 앞부분의 절이나 부록이 있다면, 귀하는 선택에 따라 이 절들의
일부 또는 전부를 변경 불가 절로 지정할 수 있습니다. 그렇게 하려면 그
제목들을 수정판의 사용 허가 고지에 있는 변경 불가 절 목록에 추가하십시오.
이 제목들은 다른 모든 절 제목과 구별되어야 합니다.

귀하는 "추천사" ("Endorsements")라는 제목의 절을 추가할 수 있습니다.
다만 그 절에는 여러 당사자가 귀하의 수정판을 추천하는 내용만 들어 있어야
합니다. 예를 들면 동료 평가를 거쳤다는 진술이나, 어느 조직이 그 텍스트를
표준의 권위 있는 정의로 승인했다는 진술이 이에 해당합니다.

귀하는 수정판의 표지 문구 목록 끝에 최대 5단어의 문구 하나를 앞표지
문구로, 최대 25단어의 문구 하나를 뒷표지 문구로 추가할 수 있습니다.
어느 한 주체든 직접 또는 그 주체가 마련한 약정을 통하여 앞표지 문구 한
개와 뒷표지 문구 한 개만 추가할 수 있습니다. 문서에 이미 같은 표지에
대한 표지 문구가 있고, 그것을 귀하가 이전에 추가했거나 귀하가 대리하는
동일한 주체가 마련한 약정에 따라 추가했다면 다른 문구를 추가할 수
없습니다. 다만 이전 문구를 추가한 이전 발행자가 명시적으로 허가하면
그 문구를 교체할 수 있습니다.

문서의 저자와 발행자는 이 사용 허가서에 의하여, 수정판을 홍보하기 위해
자신의 이름을 사용하거나 수정판을 추천한다고 주장하거나 암시하는 것을
허가하지 않습니다.


\section{문서의 결합 (COMBINING DOCUMENTS)}
\label{fdl-section-combining-documents}


귀하는 위 제4절에서 수정판에 관해 정한 조건에 따라, 문서를 이 사용
허가서에 따라 공개된 다른 문서들과 결합할 수 있습니다. 다만 결합물에
모든 원문 문서의 모든 변경 불가 절을 변경하지 않은 채 포함하고, 결합한
저작물의 사용 허가 고지에 그것들을 모두 변경 불가 절로 열거하며, 그
문서들의 모든 보증 부인 조항을 보존해야 합니다.

결합된 저작물에는 이 사용 허가서의 사본을 하나만 포함해도 되며, 서로
동일한 변경 불가 절이 여러 개이면 사본 하나로 대체할 수 있습니다. 이름은
같지만 내용이 다른 변경 불가 절이 여러 개이면, 각 절의 제목 끝에 괄호를
치고 그 절의 원 저자나 발행자의 이름을 알고 있는 경우 그 이름을, 모르는
경우 고유한 번호를 추가하여 각 제목을 고유하게 만드십시오. 결합된
저작물의 사용 허가 고지에 있는 변경 불가 절 목록의 절 제목에도 똑같이
조정하십시오.

결합물에서는 여러 원문 문서의 "이력" ("History")이라는 제목의 절들을
결합하여 "이력"이라는 제목의 절 하나로 만들어야 합니다. 마찬가지로
"감사의 글" ("Acknowledgements")이라는 제목의 절들과 "헌사"
("Dedications")라는 제목의 절들도 각각 결합해야 합니다. "추천사"
("Endorsements")라는 제목의 절은 모두 삭제해야 합니다.

\section{문서 모음 (COLLECTIONS OF DOCUMENTS)}
\label{fdl-section-collections-of-documents}

귀하는 문서와 이 사용 허가서에 따라 공개된 다른 문서들로 이루어진 모음을
만들고, 각 문서에 들어 있는 이 사용 허가서의 개별 사본들을 그 모음에
포함된 사본 하나로 대체할 수 있습니다. 다만 그 밖의 모든 면에서 각 문서를
원문 그대로 복제하는 것에 관한 이 사용 허가서의 규칙을 따라야 합니다.

귀하는 그러한 모음에서 문서 하나를 추출하여 이 사용 허가서에 따라
개별적으로 배포할 수 있습니다. 다만 추출한 문서에 이 사용 허가서의
사본을 삽입하고, 그 문서를 원문 그대로 복제하는 것에 관한 그 밖의 모든
면에서 이 사용 허가서를 따라야 합니다.


\section{독립 저작물과의 집합 (AGGREGATION WITH INDEPENDENT WORKS)}
\label{fdl-section-aggregation-with-independent-works}


문서 또는 그 파생물을 다른 별개의 독립적인 문서나 저작물과 저장 또는
배포 매체의 한 권이나 한 단위에 함께 편집한 것을 "집합물"
("aggregate")이라고 합니다. 단, 그 편집물에서 발생하는 저작권이 개별
저작물이 허용하는 범위를 넘어 편집물 이용자의 법적 권리를 제한하는 데
사용되지 않아야 합니다. 문서가 집합물에 포함되더라도, 그 집합물 속의
다른 저작물이 그 자체로 문서의 파생물이 아니라면 이 사용 허가서는 그
다른 저작물에 적용되지 않습니다.

제3절의 표지 문구 요건이 이 문서의 사본들에 적용되는 경우, 문서가 전체
집합물의 절반 미만이면 문서의 표지 문구를 집합물 안에서 문서를 앞뒤로
감싸는 표지에 둘 수 있으며, 문서가 전자 형식이면 이에 상응하는 전자적
표지에 둘 수 있습니다. 그렇지 않으면 표지 문구는 전체 집합물을 앞뒤로
감싸는 인쇄 표지에 표시되어야 합니다.


\section{번역 (TRANSLATION)}
\label{fdl-section-translation}


번역은 수정의 한 종류로 간주되므로, 귀하는 제4절의 조건에 따라 문서의
번역판을 배포할 수 있습니다. 변경 불가 절을 번역문으로 대체하려면 그
저작권자의 특별 허가가 필요하지만, 변경 불가 절의 원문 판본에 더하여
그 변경 불가 절 일부 또는 전부의 번역문을 포함할 수 있습니다. 귀하는
이 사용 허가서, 문서의 모든 사용 허가 고지 및 모든 보증 부인 조항의
번역문을 포함할 수 있습니다. 다만 이 사용 허가서의 원문 영어판과 해당
고지 및 부인 조항의 원문 판본도 함께 포함해야 합니다. 번역문과 이 사용
허가서의 원문 판본, 또는 고지나 부인 조항의 원문 판본 사이에 불일치가
있으면 원문 판본이 우선합니다.

문서의 어느 절이 "감사의 글" ("Acknowledgements"), "헌사"
("Dedications") 또는 "이력" ("History")이라는 제목의 절이라면, 그
제목을 유지해야 한다는 요구 사항(제4절)은 (제1절에 따라) 통상 실제
제목을 변경할 것을 요구합니다.


\section{종료 (TERMINATION)}
\label{fdl-section-termination}


귀하는 이 사용 허가서에서 명시적으로 정한 경우를 제외하고 문서를 복제,
수정, 재허가 또는 배포할 수 없습니다. 문서를 달리 복제, 수정, 재허가
또는 배포하려는 모든 시도는 무효이며, 이 사용 허가서에 따른
귀하의 권리를 자동으로 종료시킵니다. 다만 이 사용 허가서에 따라 귀하에게서
사본이나 권리를 받은 당사자가 이 사용 허가서를 계속 완전히 준수하는
한, 그 당사자의 사용 허가는 종료되지 않습니다.


\section{이 사용 허가서의 향후 개정판 (FUTURE REVISIONS OF THIS LICENSE)}
\label{fdl-section-future-revisions-of-this-license}


자유 소프트웨어 재단(Free Software Foundation)은 때때로 GNU 자유 문서
사용 허가서의 새 개정판을 발행할 수 있습니다. 그러한 새 판들은 현재
판과 취지가 비슷하지만, 새로운 문제나 우려 사항을 다루기 위해 세부
내용이 다를 수 있습니다. https://www.gnu.org/copyleft/ 을 참조하십시오.

이 사용 허가서의 각 판에는 서로 구별되는 판 번호가 부여됩니다. 문서에서
특정 번호의 이 사용 허가서 "또는 그 이후 판"이 적용된다고 명시하면,
귀하는 명시된 그 판 또는 자유 소프트웨어 재단이 (초안이 아닌 형태로)
발행한 그 이후 판의 조건 중 어느 하나를 따를 수 있습니다. 문서에서 이
사용 허가서의 판 번호를 명시하지 않으면, 자유 소프트웨어 재단이 지금까지
(초안이 아닌 형태로) 발행한 어느 판이든 선택할 수 있습니다.


\section{부록: 귀하의 문서에 이 사용 허가서를 사용하는 방법 (ADDENDUM: How to use this License for your documents)}
\label{fdl-section-addendum-how-to-use-this-license-for-your-documents}

귀하가 작성한 문서에 이 사용 허가서를 사용하려면, 문서에 이 사용 허가서의
사본을 포함하고 제목 페이지 바로 뒤에 다음 저작권 고지와 사용 허가 고지를
넣으십시오.

\bigskip
\begin{quote}
    저작권 \copyright  연도  귀하의 이름.
    이 문서는 자유 소프트웨어 재단이 발행한 GNU 자유 문서 사용 허가서
    버전 1.2 또는 그 이후 판의 조건에 따라 복제, 배포 및/또는 수정할 수
    있습니다. 변경 불가 절, 앞표지 문구 및 뒷표지 문구는 없습니다.
    이 사용 허가서의 사본은 "GNU 자유 문서 사용 허가서 (GNU Free
    Documentation License)"라는 제목의 절에 포함되어 있습니다.
\end{quote}
\bigskip

변경 불가 절, 앞표지 문구 및 뒷표지 문구가 있다면
"변경 불가 절, 앞표지 문구 및 뒷표지 문구는 없습니다."라는 줄을 다음으로
바꾸십시오.

\bigskip
\begin{quote}
    변경 불가 절은 [그 제목들을 나열하십시오]이고, 앞표지 문구는
    [그 문구들을 나열하십시오]이며, 뒷표지 문구는
    [그 문구들을 나열하십시오]입니다.
\end{quote}
\bigskip

표지 문구 없이 변경 불가 절만 있거나, 이 세 요소가 다른 방식으로 조합된
경우에는 상황에 맞게 위 두 대안을 합치십시오.

귀하의 문서에 프로그램 코드의 사소하지 않은 예제가 포함되어 있다면,
그 예제들을 자유 소프트웨어에서 사용할 수 있게 하기 위해 GNU 일반 공중
사용 허가서와 같이 귀하가 선택한 자유 소프트웨어 사용 허가서에 따라
병행하여 공개할 것을 권장합니다.

%---------------------------------------------------------------------
